% Fluxograma da metodologia — incluído em trabalho_final.tex por % Fluxograma da metodologia — incluído em trabalho_final.tex por % Fluxograma da metodologia — incluído em trabalho_final.tex por % Fluxograma da metodologia — incluído em trabalho_final.tex por \input{figuras/fluxo-metodologia}
% Requer no preâmbulo: \usepackage{tikz} e \usetikzlibrary{arrows.meta,positioning}
\begin{figure}[htbp]
    \centering
    \caption{Fluxo geral das etapas metodológicas}
    \label{fig:fluxo-metodologia}
    \vspace{-0.2cm}
    \begin{tikzpicture}[
        node distance=0.55cm,
        etapa/.style={
            draw=black,
            thick,
            rounded corners=2pt,
            align=center,
            text width=9.2cm,
            inner sep=6pt,
            font=\small,
            fill=white
        },
        seta/.style={-Latex, thick}
    ]
        \node[etapa] (a) {\textbf{1. Coleta e tratamento dos dados}\\ \footnotesize CNPJ, CNAE, CNEFE e malhas do IBGE integrados em banco local};
        \node[etapa, below=of a] (b) {\textbf{2. Projeto da solução}\\ \footnotesize requisitos, casos de uso, classes e protótipo das interfaces};
        \node[etapa, below=of b] (c) {\textbf{3. Implementação da solução}\\ \footnotesize programas de carga e aplicação web};
        \node[etapa, below=of c] (d) {\textbf{4. Testes}\\ \footnotesize conferência da localização, comparação com o CEMPRE e inspeção heurística};

        \draw[seta] (a) -- (b);
        \draw[seta] (b) -- (c);
        \draw[seta] (c) -- (d);
    \end{tikzpicture}
    \vspace{-0.2cm}
    \\\textbf{\footnotesize Fonte: Elaborado pelo autor.}
\end{figure}

% Requer no preâmbulo: \usepackage{tikz} e \usetikzlibrary{arrows.meta,positioning}
\begin{figure}[htbp]
    \centering
    \caption{Fluxo geral das etapas metodológicas}
    \label{fig:fluxo-metodologia}
    \vspace{-0.2cm}
    \begin{tikzpicture}[
        node distance=0.55cm,
        etapa/.style={
            draw=black,
            thick,
            rounded corners=2pt,
            align=center,
            text width=9.2cm,
            inner sep=6pt,
            font=\small,
            fill=white
        },
        seta/.style={-Latex, thick}
    ]
        \node[etapa] (a) {\textbf{1. Coleta e tratamento dos dados}\\ \footnotesize CNPJ, CNAE, CNEFE e malhas do IBGE integrados em banco local};
        \node[etapa, below=of a] (b) {\textbf{2. Projeto da solução}\\ \footnotesize requisitos, casos de uso, classes e protótipo das interfaces};
        \node[etapa, below=of b] (c) {\textbf{3. Implementação da solução}\\ \footnotesize programas de carga e aplicação web};
        \node[etapa, below=of c] (d) {\textbf{4. Testes}\\ \footnotesize conferência da localização, comparação com o CEMPRE e inspeção heurística};

        \draw[seta] (a) -- (b);
        \draw[seta] (b) -- (c);
        \draw[seta] (c) -- (d);
    \end{tikzpicture}
    \vspace{-0.2cm}
    \\\textbf{\footnotesize Fonte: Elaborado pelo autor.}
\end{figure}

% Requer no preâmbulo: \usepackage{tikz} e \usetikzlibrary{arrows.meta,positioning}
\begin{figure}[htbp]
    \centering
    \caption{Fluxo geral das etapas metodológicas}
    \label{fig:fluxo-metodologia}
    \vspace{-0.2cm}
    \begin{tikzpicture}[
        node distance=0.55cm,
        etapa/.style={
            draw=black,
            thick,
            rounded corners=2pt,
            align=center,
            text width=9.2cm,
            inner sep=6pt,
            font=\small,
            fill=white
        },
        seta/.style={-Latex, thick}
    ]
        \node[etapa] (a) {\textbf{1. Coleta e tratamento dos dados}\\ \footnotesize CNPJ, CNAE, CNEFE e malhas do IBGE integrados em banco local};
        \node[etapa, below=of a] (b) {\textbf{2. Projeto da solução}\\ \footnotesize requisitos, casos de uso, classes e protótipo das interfaces};
        \node[etapa, below=of b] (c) {\textbf{3. Implementação da solução}\\ \footnotesize programas de carga e aplicação web};
        \node[etapa, below=of c] (d) {\textbf{4. Testes}\\ \footnotesize conferência da localização, comparação com o CEMPRE e inspeção heurística};

        \draw[seta] (a) -- (b);
        \draw[seta] (b) -- (c);
        \draw[seta] (c) -- (d);
    \end{tikzpicture}
    \vspace{-0.2cm}
    \\\textbf{\footnotesize Fonte: Elaborado pelo autor.}
\end{figure}

% Requer no preâmbulo: \usepackage{tikz} e \usetikzlibrary{arrows.meta,positioning}
\begin{figure}[htbp]
    \centering
    \caption{Fluxo geral das etapas metodológicas}
    \label{fig:fluxo-metodologia}
    \vspace{-0.2cm}
    \begin{tikzpicture}[
        node distance=0.55cm,
        etapa/.style={
            draw=black,
            thick,
            rounded corners=2pt,
            align=center,
            text width=9.2cm,
            inner sep=6pt,
            font=\small,
            fill=white
        },
        seta/.style={-Latex, thick}
    ]
        \node[etapa] (a) {\textbf{1. Coleta e tratamento dos dados}\\ \footnotesize CNPJ, CNAE, CNEFE e malhas do IBGE integrados em banco local};
        \node[etapa, below=of a] (b) {\textbf{2. Projeto da solução}\\ \footnotesize requisitos, casos de uso, classes e protótipo das interfaces};
        \node[etapa, below=of b] (c) {\textbf{3. Implementação da solução}\\ \footnotesize programas de carga e aplicação web};
        \node[etapa, below=of c] (d) {\textbf{4. Testes}\\ \footnotesize conferência da localização, comparação com o CEMPRE e inspeção heurística};

        \draw[seta] (a) -- (b);
        \draw[seta] (b) -- (c);
        \draw[seta] (c) -- (d);
    \end{tikzpicture}
    \vspace{-0.2cm}
    \\\textbf{\footnotesize Fonte: Elaborado pelo autor.}
\end{figure}
